\section{Conclusion}
This case study has examined the architecture, operational principles, and 
operating system interaction of three major disk storage technologies: HDD, 
SATA SSD, and NVMe SSD. The findings confirm that while HDDs remain a cost-effective 
solution for bulk storage, their mechanical design imposes inherent limitations on 
speed, reliability, and power efficiency. SSDs address many of these limitations 
through NAND flash memory, though SATA-based SSDs are still constrained by an 
interface and command protocol originally designed for mechanical drives.

NVMe represents the most significant advancement examined in this study, as it 
was purpose-built to exploit the capabilities of flash memory and the PCIe bus. 
Through its highly parallel queue-based architecture, NVMe substantially reduces 
latency and increases throughput compared to both HDDs and SATA SSDs, making it 
the preferred choice for performance-critical and data-intensive applications.

Understanding how the operating system manages each storage technology, from file 
system creation and partitioning to specialized functions such as TRIM, wear leveling, 
and disk scheduling, is essential for making informed decisions about system design and 
storage selection. As storage demands continue to grow with the expansion of cloud computing, 
artificial intelligence, and big data applications, NVMe and its extensions, such as NVMe over 
Fabrics, are expected to play an increasingly central role in meeting the performance 
requirements of modern computing systems.